% HPC/GPU-first CV variant (labs, GPU-platform, video-infra eng roles).
% Builds to public/elia_gatti_cv_hpc.pdf.
% Shared preamble for Elia Gatti CV variants (Overleaf + local Tectonic).
% ATS-first: no icon fonts, single-column flow, standard pdfTeX hygiene.

\documentclass[a4paper,10pt]{article}

\usepackage[T1]{fontenc}
\usepackage{montserrat}
\renewcommand{\familydefault}{\sfdefault}
\usepackage{url}
\usepackage{parskip}
\RequirePackage{color}
\RequirePackage{graphicx}
\usepackage[usenames,dvipsnames]{xcolor}
\usepackage[scale=0.9, top=1.2cm, bottom=1.0cm]{geometry}
\usepackage{tabularx}
\usepackage{enumitem}
\usepackage{titlesec}
\usepackage{multicol}
\usepackage{multirow}
\usepackage[unicode, draft=false]{hyperref}
\usepackage{tikz}
\usetikzlibrary{shapes.misc}

% ATS / copy-paste hygiene (pdfTeX only; XeTeX emits Unicode natively).
\usepackage{iftex}
\ifPDFTeX
  \usepackage{cmap}
  \input{glyphtounicode}
  \pdfgentounicode=1
\fi

\definecolor{maintext}{RGB}{25,25,25}
\definecolor{accent}{RGB}{50,50,50}
\definecolor{linkcolour}{RGB}{50,50,50}
\definecolor{sectionrule}{RGB}{180,180,180}
\definecolor{dividergrey}{RGB}{210,210,210}

\titleformat{\section}
  {\Large\bfseries\color{accent}\raggedright}
  {}
  {0em}
  {}
  [\color{sectionrule}\titlerule]
\titlespacing{\section}{0pt}{5pt}{5pt}  % Tight: one-page budget

\hypersetup{colorlinks,breaklinks,urlcolor=linkcolour,linkcolor=linkcolour}

% Keep technical terms breaking at sane points only.
\hyphenation{co-a-lesc-ing}

\newcommand{\projectdivider}{%
    \vspace{1pt}
    \noindent\begin{tikzpicture}
        \draw[dividergrey, dash pattern=on 2pt off 2pt, line width=0.4pt] (0,0) -- (\linewidth,0);
    \end{tikzpicture}%
    \vspace{1pt}
}

% Single-row header (title left, date right). Body stays single-column flow.
\newenvironment{joblong}[2]{%
    \begin{tabularx}{\linewidth}{@{}l X r@{}}
    \textbf{#1} & \hfill & #2 \\[2pt]
    \end{tabularx}
    \begin{minipage}[t]{\linewidth}
    \begin{itemize}[nosep, after=\strut, leftmargin=1.5em, itemsep=2pt, label=\color{accent}\textbullet]
}{%
    \end{itemize}
    \end{minipage}%
}


\begin{document}
\pagestyle{empty}
\color{maintext}

\begin{center}
{\Huge\textbf{Elia Gatti}} \\[6pt]
{\large Software Engineer (HPC \& GPU) \ $\cdot$ \ M.S. Computer Science, University of Trento} \\[5pt]
\href{mailto:elia.gatti01@gmail.com}{elia.gatti01@gmail.com} \ $\cdot$ \
\href{https://www.linkedin.com/in/elia-gatti/}{linkedin.com/in/elia-gatti} \ $\cdot$ \
\href{https://github.com/MaiDormo}{github.com/MaiDormo} \ $\cdot$ \
\href{https://maidormo.github.io/}{maidormo.github.io} \\[3pt]
{\small Available Mar 2027}
\end{center}

\section{Work Experience}
\begin{joblong}{Bitmovin, Software Engineer Intern}{June 2026 to Sept 2026}
    \item Shipped KAIROS \textbf{end-to-end}, an AI video highlight detection and segmentation service (VOD plus live): a \textbf{public API with MCP} and Go/Python workers on Terraform-managed GCP. [\href{https://kairosapp.tech}{demo}]
    \item Cut transcription latency by \textbf{20\%} by tuning worker parallelism and slimming workers, guided by live playback measurements.
    \item Reduced failed runs and manual operations by building \textbf{job orchestration} with startup fixes and \textbf{live observability}.
    \item Backed the API and workers with a \textbf{Supabase/PostgreSQL schema}, queries, and \textbf{RLS policies}.
\end{joblong}

\vspace{4pt}

\begin{joblong}{Dedagroup, Software Engineer Intern}{May 2024 to Sept 2024}
    \item Halved VM resource requirements by migrating the `TEN' treasury application from legacy Windows to Linux servers.
    \item Delivered full-stack features with an HTMX frontend on a Java/Spring backend while maintaining the existing codebase.
\end{joblong}

\section{Education}
\textbf{M.S. in Computer Science, University of Trento} \hfill Expected Mar 2027 \\
{\small Thesis (in progress): GPU kernel optimization. GPU programming and high-performance computing.} \\[4pt]
\textbf{B.S. in Computer Science, University of Trento} \hfill 2024 \\
{\small Bachelor's thesis: MPEG-DASH performance analysis under simulated (Mininet) and real (AWS) networks.}

\section{Projects}

\textbf{Sparse Matrix-Vector Multiplication on GPU} [\href{https://github.com/MaiDormo/cuda-SpMV}{code}]
\begin{itemize}[nosep, after=\strut, leftmargin=1.5em, itemsep=2pt, label=\color{accent}\textbullet]
    \item Developed a \textbf{hybrid adaptive CUDA kernel} for NVIDIA A30, plus a V2 with sub-warp lane classes and huge-row splitting (deterministic finalize, no atomics).
    \item Verified results against a double-precision CPU reference while profiling execution time and GFLOPS (memory coalescing, occupancy, L2-cache control).
\end{itemize}

\projectdivider
\textbf{Parallel Minimum Spanning Tree (MPI + OpenMP)} [\href{https://github.com/MaiDormo/parallel_mst}{code}]
\begin{itemize}[nosep, after=\strut, leftmargin=1.5em, itemsep=2pt, label=\color{accent}\textbullet]
    \item Implemented parallel Boruvka/Kruskal MST with \textbf{MPI + OpenMP} using a flat 1D layout, mmap I/O, and bulk collectives; analyzed speedup and scalability on a CPU cluster (report in repo).
\end{itemize}

\projectdivider
\textbf{Distributed Key-Value Storage System} [\href{https://github.com/MaiDormo/distributed-storage-system}{code}]
\begin{itemize}[nosep, after=\strut, leftmargin=1.5em, itemsep=2pt, label=\color{accent}\textbullet]
    \item Built a fault-tolerant store with \textbf{Akka actors} (Java 21) using consistent hashing and quorum reads/writes, surviving concurrent access and node failures without client disruption.
\end{itemize}

\section{Achievements}
\textbf{CTM Challenge, Trento 2025 --- 1st place} [\href{https://github.com/MaiDormo/DWT-SVD-watermarking}{code}]
\begin{itemize}[nosep, after=\strut, leftmargin=1.5em, itemsep=2pt, label=\color{accent}\textbullet]
    \item Engineered SVD-based watermark embedding within DWT blocks, robust to compression, noise and transforms; validated with an automated attack suite using ROC analysis.
\end{itemize}

\projectdivider
\textbf{EuroTech HK Hackathon, Munich 2026 --- Guardian: privacy-first eldercare monitor} [\href{https://github.com/MaiDormo/Guardian}{code}]
\begin{itemize}[nosep, after=\strut, leftmargin=1.5em, itemsep=2pt, label=\color{accent}\textbullet]
    \item Built zero-config elderly monitoring with mmWave radar plus an on-device Gemma 4 LLM (no cameras, wearables, or cloud) with behavioral drift tracking and real-time crisis alerting.
\end{itemize}

\section{Skills}
\small
\textbf{Languages:} C, C++, CUDA, Java, Python, Go \\
\textbf{Tools \& Frameworks:} MPI, OpenMP, Slurm, Docker, Git, Linux, PostgreSQL, FastAPI \\
\textbf{Spoken languages:} Italian (native), English (fluent)

\end{document}
